\documentclass[border=8pt]{standalone}
\usepackage[siunitx]{circuitikz}

\begin{document}
\begin{circuitikz}[european resistors, line width=0.65pt]
  % Excitation source and complete supply/return rails.
  \draw (-2.0,0) node[ground] {}
    to[V, l_={$V_{\mathrm{EX}}$}] (-2.0,6.0)
    -- (5.5,6.0);
  \node[font=\small] at (-2.35,5.20) {$+$};
  \node[font=\small] at (-2.35,0.80) {$-$};
  \draw (-2.0,0) -- (5.5,0);

  % Four bridge arms. The left and right midpoint voltages are v+ and v-.
  \draw (1.0,6.0)
    to[R, l_={$R_1$}] (1.0,3.0)
    to[R, l_={$R_2$}] (1.0,0);
  \draw (5.5,6.0)
    to[R, l={$R_3$}] (5.5,3.0)
    to[R, l={$R_4$}] (5.5,0);

  % Test points are brought away from the bridge arms.
  \draw (1.0,3.0) -- (2.00,3.0)
    node[ocirc, label={[font=\small]above:{TP$+$: $v_+$}}] {};
  \draw (5.5,3.0) -- (6.40,3.0)
    node[ocirc, label={[font=\small]right:{TP$-$: $v_-$}}] {};
  \draw (5.5,6.0) -- (6.40,6.0)
    node[ocirc, label={[font=\small]right:{TPEX: $V_{\mathrm{EX}}$}}] {};
  \draw (5.5,0) -- (6.40,0)
    node[ocirc, label={[font=\small]right:{TP0: $0$ V}}] {};

  % Explicit differential-output polarity and definition.
  \draw (8.10,3.25) -- (8.50,3.25)
    -- (8.50,2.75) -- (8.10,2.75);
  \node[font=\small] at (7.85,3.25) {$+$};
  \node[font=\small] at (7.85,2.75) {$-$};
  \node[font=\small, align=left, anchor=west] at (8.80,3.0)
    {$v_d\equiv v_+-v_-$};

  % Only junctions of three or more conductors receive dots.
  \node[circ] at (1.0,6.0) {};
  \node[circ] at (1.0,3.0) {};
  \node[circ] at (1.0,0) {};
  \node[circ] at (5.5,6.0) {};
  \node[circ] at (5.5,3.0) {};
  \node[circ] at (5.5,0) {};

  \node[font=\small, align=center] at (3.25,-0.65)
    {excitation return and measurement reference};
\end{circuitikz}
\end{document}
